\section{Notas de Mapeamento para Firestore}
% =================================================

\begin{explicacao}[Por que esta seção existe]
A modelagem da seção anterior é deliberadamente relacional e normalizada (3FN), pensada para ser lida e avaliada como um modelo de dados independente de tecnologia. O banco real do LCQUI, porém, é o \textbf{Firestore} (Firebase), um banco de documentos sem \texttt{JOIN} e sem \texttt{LIKE}/busca por substring. Esta seção documenta as denormalizações propositais necessárias para a implementação -- nenhuma delas deve ser lida como parte do modelo 3FN, e nenhuma delas deve alterar a seção 4.
\end{explicacao}

\subsection{Reintrodução de \textit{letra\_inicial}}
O campo \textit{letra\_inicial}, removido de \textit{Resumo\_Reagente} por ser derivável de \textit{nome} (violaria 3FN), volta a existir como campo próprio \textbf{apenas} no documento Firestore da coleção \textit{Resumo\_Reagente}. Motivo: o Firestore não tem \texttt{LIKE}/busca por substring, então a busca alfabética do catálogo precisa de um campo indexável para filtro de igualdade (ou de \textit{range query} por prefixo). Como esse filtro normalmente é combinado com outros filtros de igualdade na mesma consulta (estado físico, natureza química), um campo de igualdade simples é mais previsível dentro das limitações de índice composto do Firestore do que uma \textit{range query} de prefixo sobre \textit{nome} combinada com filtros adicionais.

\begin{regra}[Consistência obrigatória]
\textit{letra\_inicial} deve ser recalculado e reescrito pela mesma \textit{Cloud Function} que grava \textit{nome}, nunca aceito diretamente do cliente -- ver princípio ``nunca confiar no \textit{client-side}'' na seção \textit{Tecnologia e Relatórios}.
\end{regra}

\subsection{Estado físico canônico do catálogo}
DP-A01: Resumo\_Reagente.estado\_fisico é escalar SOLIDO/LIQUIDO; eh\_higroscopico também pertence ao resumo. Filtro usa igualdade: \texttt{.where("estado\_fisico", "==", "SOLIDO")}. A unidade operacional é derivada, g para sólido e mL para líquido.

\subsection{Espelhamento de \textit{nome\_equipamento} e Localização em \textit{Bem\_Patrimonial}}
No modelo 3FN, o nome do tipo de bem pertence exclusivamente a \textit{Resumo\_Bem\_Patrimonial.nome}; \textit{Bem\_Patrimonial} mantém apenas \textit{id\_resumo\_bem\_patrimonial} e \textit{id\_local}. No Firestore, porém, consultas de listagem e relatórios frequentemente partem diretamente de \textit{Bem\_Patrimonial}, sem um \textit{JOIN} para resolver o resumo ou o prédio. Por isso, o documento Firestore de \textit{Bem\_Patrimonial} deve conter os campos denormalizados \textit{nome\_equipamento}, \textit{predio}, \textit{andar} e \textit{sala}.

\begin{regra}[Fonte única de verdade para o nome do patrimônio]
\textit{nome\_equipamento} é exclusivamente um campo de leitura derivado de \textit{Resumo\_Bem\_Patrimonial.nome}. O cliente nunca pode gravá-lo diretamente. Sempre que o nome do resumo for alterado, a mesma operação transacional ou um gatilho de propagação controlado pelo servidor deve atualizar os documentos \textit{Bem\_Patrimonial} associados. Na aprovação da requisição de edição descrita na seção 10.2.5, \textit{novo\_nome} representa uma proposta de alteração do nome do resumo; o servidor resolve/cria outro resumo e altera a referência somente do bem solicitado; renomear globalmente um resumo é ação distinta do gestor.
\end{regra}

\subsection{Composição embutida na especificação}
O mapeamento físico adotado é \texttt{composicao}, array de mapas na especificação em \texttt{Resumo\_Reagente/id/Especificacoes/id}. Cada mapa referencia uma substância; o servidor impede IDs repetidos e limita o tamanho do array. Não manter simultaneamente uma coleção raiz e uma subcoleção alternativa como fontes concorrentes. A relação N:N continua normalizada na seção 4.

\subsection{Tabelas ``Materialized'' são coleções próprias, não \textit{views} do SGBD}
As tabelas descritas na seção \textit{Materialized (Views)} não são \textit{materialized views} no sentido de PostgreSQL (não existe \texttt{REFRESH MATERIALIZED VIEW} no Firestore). Cada uma delas é implementada como sua própria coleção Firestore, escrita por \textit{Cloud Functions} -- gatilhos de escrita (\texttt{onDocumentWritten}) para contadores que reagem a um evento pontual (ex.: \textit{Lote\_Materializado}), ou funções agendadas (\texttt{onSchedule}) para os resumos diários/mensais. Ver pseudo-código na seção \textit{Tecnologia e Relatórios}.

\subsection{Unicidade de \texttt{numero\_patrimonio\_proposto} no Firestore}
O PostgreSQL implementa a regra RF10b via índice único parcial
\texttt{UNIQUE(numero\_patrimonio\_proposto) WHERE status = 'pendente'},
impedindo duas requisições simultâneas para a mesma plaqueta enquanto
o bem patrimonial ainda não existe. No Firestore não há índice parcial
nem \emph{unique constraint} nativo, então a garantia é reimplementada
na Cloud Function de criação da requisição de adição.

A criação lê o documento determinístico em \texttt{Locks\_Requisicao\_Patrimonio}, verifica a plaqueta já cadastrada e grava lock e requisição na mesma transação. Aprovação e rejeição removem o lock na mesma transação da resposta. Uma consulta prévia pode melhorar a mensagem de erro, mas não substitui a chave de unicidade. Locks vinculados a requisições pendentes não expiram por idade.

\subsection{Mecanismo de Sequenciamento Atômico via Documento Singleton (\textit{Contador\_Codigo\_Frasco})}
A coleção \texttt{Contador\_Codigo\_Frasco} não existe no modelo 3FN porque bancos relacionais possuem primitivas de \texttt{SEQUENCE/SERIAL}. Para gerar o identificador atômico \texttt{LCQUI-X}, utiliza-se o documento singleton \texttt{Contador\_Codigo\_Frasco/singleton} com o campo \texttt{ultimo\_codigo\_gerado}, operado atomicamente via \texttt{runTransaction} no backend. A transação garante a unicidade; contenção e custo devem ser medidos incluindo histórico, auditoria, índices e tentativas repetidas. Não se fixa aqui uma cota universal de escritas por segundo.

\subsection{Índices Compostos Obrigatórios}
Para garantir a integridade das \textit{queries} e ordenações cruciais, os seguintes índices compostos devem ser criados explicitamente no console do Firestore:
\begin{enumerate}
    \item \texttt{Historico\_Frasco\_Reagente}: \texttt{id\_frasco\_reagente (ASC)} + \texttt{timestamp (DESC)}.
    \item \texttt{Emprestimo\_Reagente}: \texttt{id\_frasco\_reagente (ASC)} + \texttt{data\_retirada (DESC)}.
    \item \textbf{M-23} \texttt{Bem\_Patrimonial}: \texttt{predio (ASC)} + \texttt{andar (ASC)} + \texttt{sala (ASC)} -- suporta o relatório personalizado por localização (Seção~9 fluxo ALM-01).
    \item \textbf{M-23} \texttt{Bem\_Patrimonial}: \texttt{id\_resumo\_bem\_patrimonial (ASC)} + \texttt{status (ASC)} -- suporta a listagem de bens por tipo e status.
    \item \textbf{M-23 / N-08} Collection-group \texttt{Historico} (sub-coleção de \texttt{Bem\_Patrimonial}): \texttt{predio (ASC)} + \texttt{timestamp (DESC)} -- suporta relatório histórico por prédio/período.
    \item \textbf{M-23} Collection-group \texttt{Historico}: \texttt{id\_gestor (ASC)} + \texttt{timestamp (DESC)} -- suporta relatório de atividade por operador.
    \item \textbf{M-23} \texttt{Emprestimo\_Reagente}: \texttt{id\_almoxarifado (ASC)} + \texttt{data\_retirada (DESC)} -- suporta relatório mensal por almoxarifado.
    \item \texttt{Frasco\_Reagente}: \texttt{id\_almoxarifado (ASC)} + \texttt{id\_resumo\_reagente (ASC)} + \texttt{id\_especificacao\_reagente (ASC)} + \texttt{disponibilidade (ASC)} + \texttt{estado\_fisico\_frasco (ASC)} + \texttt{em\_quarentena (ASC)} + \texttt{vencido (ASC)} -- suporta job agendado de escassez.
\end{enumerate}
\subsection{Mapeamento Completo de Coleções e Sub-coleções Firestore}
\label{subsec:mapeamento-colecoes}

\begin{explicacao}[Como ler esta tabela]
Cada linha descreve uma entidade 3FN da Seção 4, a decisão de estrutura no Firestore (coleção raiz, sub-coleção ou array de maps embutido) e a justificativa. Campos de consulta necessários podem ser projetados no documento Firestore conforme o dicionário; o símbolo \textbf{Denorm:} lista esses campos. A consistência usa transação para invariantes ou propagação versionada para apresentação; snapshots históricos não são atualizados com o cadastro atual.
\end{explicacao}

\small

\small
\setlength{\tabcolsep}{4.5pt}
\small
\setlength{\tabcolsep}{4pt}
\begin{longtable}{>{\raggedright\arraybackslash}p{0.28\textwidth} >{\raggedright\arraybackslash}p{0.26\textwidth} >{\raggedright\arraybackslash}p{0.38\textwidth}}
\toprule
\textbf{Entidade 3FN} & \textbf{Estrutura Firestore} & \textbf{Justificativa / Denorm}\\
\midrule
\endfirsthead
\toprule
\textbf{Entidade 3FN} & \textbf{Estrutura Firestore} & \textbf{Justificativa / Denorm} (continuação)\\
\midrule
\endhead
\bottomrule
\endfoot
Usuario & Coleção raiz \texttt{Usuarios} (docId = Firebase Auth UID) & Acesso direto por auth; estrutura natural do SDK.\\
Chefe\_Geral & Coleção raiz \texttt{Chefe\_Geral} (docId = uid) & Lido por \texttt{atualizarCustomClaims(uid)} para compor o claim \texttt{roles} do token JWT. Claims são projeção; conta ativa e versão de permissões devem ser verificadas conforme seção 11.\\
Gestor\_Almoxarifado & Coleção raiz \texttt{Gestor\_Almoxarifado} (docId = uid) & Idem acima.\\
Gestor\_Bens\_Patrimoniais & Coleção raiz \texttt{Gestor\_Bens\_Patrimoniais} (docId = uid) & Idem.\\
Professor & Coleção raiz \texttt{Professor} (docId = uid) & Idem.\\
Aluno & Coleção raiz \texttt{Aluno} (docId = uid). Denorm: \texttt{letra\_inicial} & Idem. \texttt{letra\_inicial} permite filtro de igualdade na busca.\\
Bolsista & Coleção raiz \texttt{Bolsista} (docId = uid) & Idem.\\
Gestor\_Almoxarifado\_x\_Almoxarifado & Coleção raiz & Query bidirecional: ``gestores de X'' e ``almoxarifados de Y''.\\
Almoxarifado & Coleção raiz & Poucas dezenas de documentos.\\
Local & Coleção raiz & Referenciado por múltiplas entidades; campos \texttt{predio}, \texttt{andar}, \texttt{sala} denormalizados em \texttt{Bem\_Patrimonial} via trigger.\\
Resumo\_Bem\_Patrimonial & Coleção raiz. Denorm: \texttt{letra\_inicial\_nome} & Agrupador de bens; filtro alfabético por letra inicial.\\
Bem\_Patrimonial & Coleção raiz. Denorm: \texttt{nome\_equipamento}, \texttt{predio}, \texttt{andar}, \texttt{sala}, \texttt{letra\_inicial\_nome} & Entidade central pesquisada diretamente; evita JOINs em listagens e relatórios.\\
Historico\_Bem\_Patrimonial & Sub-coleção \nolinkurl{Bem\_Patrimonial/\{id\}/Historico}. Denorm (N-08): \texttt{predio}, \texttt{andar}, \texttt{sala} -- snapshot imutável do local no momento do evento & Sempre consultado a partir de um bem conhecido. O snapshot de localização é gravado no momento da criação do evento histórico e não é atualizado quando o bem muda de sala; isso garante rastreabilidade do bem antes e depois da mudança. Requerido pelo relatório personalizado por prédio/período.\\
Alteracao\_Bem\_Patrimonial & Campo \texttt{alteracoes} (array de maps) dentro de cada doc de \texttt{Historico} & Poucos campos por alteração; relação 1:N estreita, nunca consultada independentemente.\\
Requisicao\_Edicao\_Bem\_Patrimonial & Coleção raiz. Denorm: \texttt{nome\_equipamento}, \texttt{numero\_patrimonio} & Listada independentemente nos painéis de gestores por qualquer filtro de status.\\
Requisicao\_Adicao\_Bem\_Patrimonial & Coleção raiz & Mesma lógica; query por status e número de patrimônio proposto.\\
Impressao\_Etiqueta\_Frasco & Coleção raiz & Refere-se a impressão em lote (virgens). Reimpressões avulsas de frascos já cadastrados usam o Registro de Auditoria.\\
Contador\_Codigo\_Frasco & Documento singleton em \nolinkurl{Contador\_Codigo\_Frasco/singleton} & Sequenciador atômico operado via \texttt{runTransaction} para gerar o \texttt{LCQUI-X}.\\
Substancia\_Quimica & Coleção raiz & Catálogo global; raramente alterado.\\
Resumo\_Reagente & Coleção raiz. Denorm: \texttt{letra\_inicial} & Catálogo de reagentes; filtros de igualdade combinados exigem ambos os campos como propriedades do documento.\\
Especificacao\_Reagente & Sub-coleção \nolinkurl{Resumo\_Reagente/\{id\}/Especificacoes} & Sempre acessada a partir do resumo; nunca consultada globalmente.\\
Composicao\_Reagente & Array de maps embutido no documento de \texttt{Especificacao\_Reagente} & N pequeno e fixo (poucas substâncias por especificação); não cresce sem limite.\\
Lote & Coleção raiz. Denorm: \texttt{id\_especificacao\_reagente}, \texttt{nome\_reagente} & Relatórios consultam todos os lotes de um período sem partir de uma especificação conhecida.\\
Frasco\_Reagente & Coleção raiz. Denorm: \texttt{id\_especificacao\_reagente}, \texttt{id\_resumo\_reagente}, \texttt{nome\_reagente}, \texttt{unidade\_medida}, \texttt{id\_almoxarifado} & Entidade central; queries cruzadas constantes por almoxarifado, status e vencimento.\\
Historico\_Frasco\_Reagente & Coleção raiz (não sub-coleção). Denorm: \texttt{id\_almoxarifado} & Jobs diários consultam ``todos os históricos do dia'' sem partir de um frasco conhecido; coleção raiz simplifica esse padrão; collection-group também permitiria históricos em subcoleções.\\
Emprestimo\_Reagente & Coleção raiz. Denorm: \texttt{id\_almoxarifado}, \texttt{nome\_reagente}, \texttt{nome\_usuario\_retirou} & Relatórios por almoxarifado e por período; queries de status global (\texttt{EM\_USO}).\\
Materia & Coleção raiz & Catálogo global; raramente alterado.\\
Professor\_x\_Materia & Coleção raiz & Query bidirecional: ``matérias de X'' e ``professores de Y''.\\
Turma & Coleção raiz. Denorm: \texttt{nome\_materia} & Acesso direto por \texttt{codigo\_turma}; nome da matéria necessário na listagem sem JOIN.\\
\textbf{\texttt{Aluno\_x\_Turma}} & \textbf{Subcoleção dupla (espelhada):}\newline 1. \nolinkurl{Turma/\{id\}/Alunos}\newline 2. \nolinkurl{Usuarios/\{uid\}/Turmas} & \textbf{1. \nolinkurl{Turma/\{id\}/Alunos}:} Visão da disciplina; contagem de capacidade e listagem de chamada para o professor.\newline \textbf{2. \nolinkurl{Usuarios/\{uid\}/Turmas}:} Visão do aluno; viabiliza o feed em tempo real no menu lateral de ``Minhas Turmas'' com segurança isolada por UID e zero leituras $N+1$. Consistência garantida via transação nas Cloud Functions de matrícula/exclusão.\\
Historico\_Alunos\_Turma & Sub-coleção \nolinkurl{Turma/\{id\}/HistoricoAlunos} & Auditoria de inclusão/exclusão sempre por turma.\\
Post & Sub-coleção \nolinkurl{Turma/\{id\}/Posts} & Feed cronológico da turma; nunca consultado globalmente.\\
Comentario & Sub-coleção \nolinkurl{Turma/\{id\}/Posts/\{postId\}/Comentarios} & Thread aninhada ao post.\\
Historico\_Posts\_Turma & Sub-coleção \nolinkurl{Turma/\{id\}/Posts/\{postId\}/Historico} & Auditoria de edição sempre a partir do post.\\
Historico\_Comentario & Sub-coleção \nolinkurl{Turma/\{id\}/Posts/\{postId\}/Comentarios/\{id\}/Historico} & Auditoria de edição sempre a partir do comentário.\\
Roteiro\_Experimento & Coleção raiz & Compartilhamento cross-professor exige acesso sem partir de uma turma.\\
Convite\_Aluno & Coleção raiz & Query por e-mail normalizado e status independentemente de turma.\\
Notificacao & Sub-coleção \nolinkurl{Usuarios/\{uid\}/Notificacoes} & Query natural por usuário; funciona automaticamente para usuários multi-role sem necessidade de multiple reads.\\
Registro\_de\_Auditoria & Coleção raiz & Log global queryável por tipo de entidade e período.\\
Lote\_Materializado & Coleção raiz (docId = id do lote) & Atualizada por trigger \texttt{onFrascoCriado}/\texttt{onFrascoRemovido}.\\
Resumo\_Almoxarifado\_Diario & Coleção raiz & Job agendado diário.\\
Resumo\_Reagente\_Diario & Coleção raiz & Job agendado diário.\\
Resumo\_Bem\_Patrimonial\_Diario & Coleção raiz & Job agendado diário.\\
Atividade\_Gestor\_Almoxarifado\_Mensal & Coleção raiz & Job agendado diário (acumulado por mês).\\
Atividade\_Gestor\_Bens\_Patrimoniais\_Mensal & Coleção raiz & Job agendado diário (acumulado por mês).\\
\end{longtable}


\begin{regra}[Regra geral de denormalização]
Copiar apenas atributos necessários às consultas documentadas. Cada projeção identifica sua origem e mecanismo de atualização; snapshots de eventos são imutáveis. Vínculos e invariantes críticos são transacionais, enquanto projeções de apresentação podem ser assíncronas com reconciliação.
\end{regra}


\normalsize
\begingroup
\sloppy
\subsubsection{Dicionário de dados completo e convenções físicas}
\label{sec:dicionario}
Este dicionário especifica o alvo documental. Divergências com os caminhos do código atual são listadas na auditoria; não se pressupõe migração executada. IDs relacionais SERIAL tornam-se docIds string; FKs tornam-se IDs string, sem JOIN implícito. Datas civis usam string YYYY-MM-DD; instantes usam Timestamp UTC, exibidos em America/Sao\_Paulo. Valores quantitativos usam number finito, com precisão/escala da seção 4; rejeitar NaN, infinito, sinal e casas decimais incompatíveis. Todos os limites VARCHAR continuam válidos. O = obrigatório; N = null permitido; C = condicional. Campos calculados, autoria, timestamps, IDs e estados são gravados pelo servidor. Arrays de composição e alterações são limitados pelo esquema, nunca históricos crescentes. O docId substitui o campo id SERIAL; em associações a chave é determinística pelas FKs.

\paragraph{Usuario}\mbox{}\par
Caminho: \texttt{Usuarios/\{uid\}}.
\begin{description}
\item[\texttt{id}] docId string; O. Chave do documento, gerada pelo servidor; não duplicar como contador relacional.
\item[\texttt{nome}] string; O. Nome de apresentação; não substitui a chave do documento. Máximo 150 caracteres.
\item[\texttt{email}] string; O. Endereço normalizado para identidade/convite; não divulgar em catálogo público interno. Máximo 150 caracteres. Unicidade garantida por chave determinística/registro de unicidade no servidor.
\item[\texttt{profile\_url\_photo}] string; N. Foto opcional de apresentação do usuário; ausência usa avatar padrão, com acesso validado ao objeto.
\end{description}

\paragraph{Chefe\_Geral}\mbox{}\par
Caminho: \texttt{Chefe\_Geral/\{id\}}.
\begin{description}
\item[\texttt{id\_usuario}] string; O. Identificador da entidade usuario, cuja existência e escopo são verificados no servidor.
\end{description}

\paragraph{Gestor\_Almoxarifado}\mbox{}\par
Caminho: \texttt{Gestor\_Almoxarifado/\{id\}}.
\begin{description}
\item[\texttt{id\_usuario}] string; O. Identificador da entidade usuario, cuja existência e escopo são verificados no servidor.
\end{description}

\paragraph{Gestor\_Almoxarifado\_x\_Almoxarifado}\mbox{}\par
Caminho: \texttt{Gestor\_Almoxarifado\_x\_Almoxarifado/\{id\}}.
\begin{description}
\item[\texttt{id\_gestor\_almoxarifado}] string; O. Identificador da entidade gestor almoxarifado, cuja existência e escopo são verificados no servidor.
\item[\texttt{id\_almoxarifado}] string; O. Identificador da entidade almoxarifado, cuja existência e escopo são verificados no servidor.
\end{description}

\paragraph{Gestor\_Bens\_Patrimoniais}\mbox{}\par
Caminho: \texttt{Gestor\_Bens\_Patrimoniais/\{id\}}.
\begin{description}
\item[\texttt{id\_usuario}] string; O. Identificador da entidade usuario, cuja existência e escopo são verificados no servidor.
\end{description}

\paragraph{Bem\_Patrimonial}\mbox{}\par
Caminho: \texttt{Bem\_Patrimonial/\{id\}}.
\begin{description}
\item[\texttt{id}] docId string; O. Chave do documento, gerada pelo servidor; não duplicar como contador relacional.
\item[\texttt{id\_resumo\_bem\_patrimonial}] string; O. Identificador da entidade resumo bem patrimonial, cuja existência e escopo são verificados no servidor.
\item[\texttt{numero\_patrimonio}] string; O. Plaqueta institucional em texto, preservando zeros iniciais; identifica unicamente a unidade física. Máximo 30 caracteres.
\item[\texttt{estado\_conservacao}] string enum; O. Valor de estado conservacao, validado pelo formulário e pelo servidor conforme regras do domínio. Valores: BOM, REGULAR, RUIM.
\item[\texttt{id\_local}] string; O. Identificador da entidade local, cuja existência e escopo são verificados no servidor.
\item[\texttt{photo\_url}] string; O. Referência da foto validada; não confiar em URL arbitrária.
\item[\texttt{documento\_dado\_baixa\_pdf\_url}] string; N. Referência do comprovante PDF, nula antes da baixa e obrigatória em Ja\_dado\_baixa; preservar após vínculo.
\item[\texttt{nome\_responsavel\_sei}] string; O. Responsável institucional em texto; não é FK de usuário. Máximo 150 caracteres.
\item[\texttt{status}] string enum; O. Estado controlado pelo domínio; transições são validadas no servidor. Valores: Ativo, Inservivel, Ja\_dado\_baixa.
\item[\texttt{descricao\_complementar}] string; N. Observação opcional da unidade física, sem alterar a descrição compartilhada do modelo. Máximo 200 caracteres.
\end{description}

\paragraph{Resumo\_Bem\_Patrimonial}\mbox{}\par
Caminho: \texttt{Resumo\_Bem\_Patrimonial/\{id\}}.
\begin{description}
\item[\texttt{id}] docId string; O. Chave do documento, gerada pelo servidor; não duplicar como contador relacional.
\item[\texttt{nome}] string; O. Nome de apresentação; não substitui a chave do documento. Máximo 150 caracteres.
\item[\texttt{descricao}] string; O. Valor de descricao, validado pelo formulário e pelo servidor conforme regras do domínio.
\end{description}

\paragraph{Historico\_Bem\_Patrimonial}\mbox{}\par
Caminho: \texttt{Bem\_Patrimonial/\{id\}/Historico/\{eventoId\}}.
\begin{description}
\item[\texttt{id}] docId string; O. Chave do documento, gerada pelo servidor; não duplicar como contador relacional.
\item[\texttt{id\_bem\_patrimonial}] string; O. Identificador da entidade bem patrimonial, cuja existência e escopo são verificados no servidor.
\item[\texttt{timestamp}] Timestamp; O. Data/instante de timestamp, com ordem temporal validada e autoria do relógio do servidor para eventos.
\item[\texttt{tipo}] string enum; O. Valor de tipo, validado pelo formulário e pelo servidor conforme regras do domínio. Valores: cadastro, edicao, baixa.
\item[\texttt{id\_usuario}] string; O. Identificador da entidade usuario, cuja existência e escopo são verificados no servidor.
\end{description}

\paragraph{Alteracao\_Bem\_Patrimonial}\mbox{}\par
Caminho: \texttt{Bem\_Patrimonial/\{id\}/Historico/\{eventoId\}: alteracoes[]}.
\begin{description}
\item[\texttt{campo}] string; O. Valor de campo, validado pelo formulário e pelo servidor conforme regras do domínio. Máximo 60 caracteres.
\item[\texttt{valor\_anterior}] string; O. Valor de valor anterior, validado pelo formulário e pelo servidor conforme regras do domínio.
\item[\texttt{valor\_novo}] string; O. Valor de valor novo, validado pelo formulário e pelo servidor conforme regras do domínio.
\end{description}

\paragraph{Requisicao\_Edicao\_Bem\_Patrimonial}\mbox{}\par
Caminho: \texttt{Requisicao\_Edicao\_Bem\_Patrimonial/\{id\}}.
\begin{description}
\item[\texttt{id}] docId string; O. Chave do documento, gerada pelo servidor; não duplicar como contador relacional.
\item[\texttt{id\_bem\_patrimonial}] string; O. Identificador da entidade bem patrimonial, cuja existência e escopo são verificados no servidor.
\item[\texttt{status}] string enum; O. Estado controlado pelo domínio; transições são validadas no servidor. Valores: pendente, aprovada, rejeitada.
\item[\texttt{feita\_em}] Timestamp; O. Data/instante de feita em, com ordem temporal validada e autoria do relógio do servidor para eventos.
\item[\texttt{respondida\_em}] Timestamp; N. Data/instante de respondida em, com ordem temporal validada e autoria do relógio do servidor para eventos.
\item[\texttt{id\_usuario\_solicitante}] string; O. Identificador da entidade usuario solicitante, cuja existência e escopo são verificados no servidor.
\item[\texttt{id\_usuario\_respondente}] string; N. Identificador da entidade usuario respondente, cuja existência e escopo são verificados no servidor.
\item[\texttt{novo\_nome}] string; N. Valor novo nome, restrito ao campo correspondente da entidade de origem; preservar ausência legítima. Máximo 150 caracteres.
\item[\texttt{novo\_status}] string enum; N. Valor novo status, restrito ao campo correspondente da entidade de origem; preservar ausência legítima. Valores: Ativo, Inservivel, Ja\_dado\_baixa.
\item[\texttt{novo\_estado\_conservacao}] string enum; N. Valor novo estado conservacao, restrito ao campo correspondente da entidade de origem; preservar ausência legítima. Valores: BOM, REGULAR, RUIM.
\item[\texttt{novo\_id\_local}] string; N. Valor novo id local, restrito ao campo correspondente da entidade de origem; preservar ausência legítima.
\item[\texttt{nova\_photo\_url}] string; N. Valor nova photo url, restrito ao campo correspondente da entidade de origem; preservar ausência legítima.
\item[\texttt{motivo}] string; O. Valor de motivo, validado pelo formulário e pelo servidor conforme regras do domínio.
\item[\texttt{justificativa\_resposta}] string; N. Justificativa obrigatória ao aprovar ou rejeitar; ausente enquanto pendente.
\end{description}

\paragraph{Requisicao\_Adicao\_Bem\_Patrimonial}\mbox{}\par
Caminho: \texttt{Requisicao\_Adicao\_Bem\_Patrimonial/\{id\}}.
\begin{description}
\item[\texttt{id}] docId string; O. Chave do documento, gerada pelo servidor; não duplicar como contador relacional.
\item[\texttt{numero\_patrimonio\_proposto}] string; O. Plaqueta solicitada para o novo bem; unicidade entre pendências e bens existentes. Máximo 30 caracteres.
\item[\texttt{status}] string enum; O. Estado controlado pelo domínio; transições são validadas no servidor. Valores: pendente, aprovada, rejeitada.
\item[\texttt{feita\_em}] Timestamp; O. Data/instante de feita em, com ordem temporal validada e autoria do relógio do servidor para eventos.
\item[\texttt{id\_bem\_patrimonial\_se\_aprovado}] string; N. Identificador da entidade bem patrimonial se aprovado, cuja existência e escopo são verificados no servidor.
\item[\texttt{respondida\_em}] Timestamp; N. Data/instante de respondida em, com ordem temporal validada e autoria do relógio do servidor para eventos.
\item[\texttt{id\_usuario\_solicitante}] string; O. Identificador da entidade usuario solicitante, cuja existência e escopo são verificados no servidor.
\item[\texttt{id\_usuario\_respondente}] string; N. Identificador da entidade usuario respondente, cuja existência e escopo são verificados no servidor.
\item[\texttt{estado\_conservacao\_proposto}] string enum; O. Valor de estado conservacao proposto, validado pelo formulário e pelo servidor conforme regras do domínio. Valores: BOM, REGULAR, RUIM.
\item[\texttt{photo\_url\_proposta}] string; O. Foto obrigatória na submissão, validada pelo backend. Valor de photo url proposta, validado pelo formulário e pelo servidor conforme regras do domínio.
\item[\texttt{nome\_responsavel\_proposto}] string; O. Valor de nome responsavel proposto, validado pelo formulário e pelo servidor conforme regras do domínio. Máximo 150 caracteres.
\item[\texttt{id\_local}] string; O. Identificador da entidade local, cuja existência e escopo são verificados no servidor.
\item[\texttt{id\_resumo\_bem\_patrimonial}] string; N. Identificador da entidade resumo bem patrimonial, cuja existência e escopo são verificados no servidor.
\item[\texttt{nome\_resumo\_proposto}] string; N. Valor de nome resumo proposto, validado pelo formulário e pelo servidor conforme regras do domínio. Máximo 150 caracteres.
\item[\texttt{descricao\_resumo\_proposta}] string; N. Valor de descricao resumo proposta, validado pelo formulário e pelo servidor conforme regras do domínio.
\item[\texttt{motivo}] string; O. Valor de motivo, validado pelo formulário e pelo servidor conforme regras do domínio.
\end{description}

\paragraph{Locks\_Requisicao\_Patrimonio}\mbox{}\par
Caminho: \texttt{Locks\_Requisicao\_Patrimonio/\{id\}}.
\begin{description}
\item[\texttt{id (docId determinístico)}] string; O. Identificador da entidade id (docId determinístico), cuja existência e escopo são verificados no servidor. Máximo 100 caracteres.
\item[\texttt{criado\_em}] Timestamp; O. Data/instante de criado em, com ordem temporal validada e autoria do relógio do servidor para eventos.
\end{description}

\paragraph{Impressao\_Etiqueta\_Frasco}\mbox{}\par
Caminho: \texttt{Impressao\_Etiqueta\_Frasco/\{id\}}.
\begin{description}
\item[\texttt{id}] docId string; O. Chave do documento, gerada pelo servidor; não duplicar como contador relacional.
\item[\texttt{gerado\_em}] Timestamp; O. Data/instante de gerado em, com ordem temporal validada e autoria do relógio do servidor para eventos.
\item[\texttt{gerado\_por}] string; O. UID do operador que solicitou a geração de etiquetas, obtido da sessão autenticada.
\item[\texttt{codigo\_inicial}] number; O. Primeiro número inteiro positivo do intervalo inclusivo de etiquetas virgens.
\item[\texttt{codigo\_final}] number; O. Último número do intervalo inclusivo, maior ou igual ao inicial e limitado a 50 etiquetas.
\end{description}

\paragraph{Almoxarifado}\mbox{}\par
Caminho: \texttt{Almoxarifado/\{id\}}.
\begin{description}
\item[\texttt{id}] docId string; O. Chave do documento, gerada pelo servidor; não duplicar como contador relacional.
\item[\texttt{id\_local}] string; O. Identificador da entidade local, cuja existência e escopo são verificados no servidor.
\item[\texttt{nome}] string; O. Nome de apresentação; não substitui a chave do documento. Máximo 100 caracteres.
\item[\texttt{descricao}] string; O. Valor de descricao, validado pelo formulário e pelo servidor conforme regras do domínio. Máximo 500 caracteres.
\end{description}

\paragraph{Estoques\_Configurados}\mbox{}\par
Caminho: \texttt{Almoxarifado/\{almoxId\}/Estoques\_Configurados/\{configId\}}.
\begin{description}
\item[\texttt{configId}] docId string; O. Chave do documento, gerada de forma determinística: \texttt{\{id\_resumo\_reagente\}\_\{id\_especificacao\_reagente\}}, garantindo unicidade do par.
\item[\texttt{id\_especificacao\_reagente}] string; O. Identificador da especificação do reagente configurada.
\item[\texttt{id\_resumo\_reagente}] string; O. Identificador do resumo do reagente, mantido no documento para consultas otimizadas e manter a identidade composta.
\item[\texttt{id\_almoxarifado}] string; O. Identificador do almoxarifado (redundante à URL, porém canônico).
\item[\texttt{qtd\_limiar\_escassez}] number; O. Limiar inteiro não negativo de quantidade de frascos para este almoxarifado.
\item[\texttt{ativo}] boolean; O. Indica se o par continua configurado/intencionalmente estocado.
\item[\texttt{notificacao\_ativa}] boolean; O. Indica se a emissão automática do alerta está habilitada (permitindo silenciar sem apagar a configuração).
\end{description}

\paragraph{Substancia\_Quimica}\mbox{}\par
Caminho: \texttt{Substancia\_Quimica/\{id\}}.
\begin{description}
\item[\texttt{id}] docId string; O. Chave do documento, gerada pelo servidor; não duplicar como contador relacional.
\item[\texttt{nome}] string; O. Nome de apresentação; não substitui a chave do documento. Máximo 150 caracteres.
\item[\texttt{cas\_number}] string; N. Identificador CAS opcional da substância; quando informado, validar formato, dígito e unicidade. Máximo 15 caracteres.
\item[\texttt{formula\_quimica}] string; N. Fórmula declarada da substância; apresentação textual, sem executar marcação arbitrária. Máximo 50 caracteres.
\item[\texttt{ativo}] boolean; O. Habilitação lógica; desativação preserva identidade e histórico.
\end{description}

\paragraph{Resumo\_Reagente}\mbox{}\par
Caminho: \texttt{Resumo\_Reagente/\{id\}}.
\begin{description}
\item[\texttt{id}] docId string; O. Chave do documento, gerada pelo servidor; não duplicar como contador relacional.
\item[\texttt{nome}] string; O. Nome de apresentação; não substitui a chave do documento. Máximo 150 caracteres.
\item[\texttt{tipo\_substancia}] string enum; O. Valor de tipo substancia, validado pelo formulário e pelo servidor conforme regras do domínio. Valores: PURA, MISTURA.
\item[\texttt{natureza\_quimica}] string enum; O. Classificação curada do item do catálogo, independente de fabricante ou frasco. Valores: ORGANICO, INORGANICO, ELEMENTO, HIBRIDO.
\item[\texttt{requer\_pesagem\_frequente}] boolean; O. Indica necessidade de alertas periódicos de pesagem; true exige frequência positiva.
\item[\texttt{frequencia\_pesagem\_dias}] number; N. Intervalo inteiro positivo entre verificações; null quando não exige pesagem frequente.
\item[\texttt{estado\_fisico}] string enum; O. SOLIDO/LIQUIDO; fonte canônica do catálogo, filtro por igualdade.
\item[\texttt{eh\_higroscopico}] boolean; O, padrão false. Fonte canônica da higroscopicidade; não propagar mudanças para snapshots de frascos já cadastrados.
\end{description}

\paragraph{Composicao\_Reagente}\mbox{}\par
Caminho: \texttt{Resumo\_Reagente/\{resumoId\}/Especificacoes/\{id\}: composicao[]}.
\begin{description}
\item[\texttt{id\_substancia\_quimica}] string; O. Identificador da entidade substancia quimica, cuja existência e escopo são verificados no servidor.
\item[\texttt{valor\_min}] number; N. Concentração pontual ou limite inferior não negativo.
\item[\texttt{valor\_max}] number; N. Limite superior; quando informado exige mínimo e valor\_min <= valor\_max.
\item[\texttt{tipo\_concentracao}] string enum; O. Base da concentração; enum M\_M, V\_V, M\_V, MOL\_L, MOL\_KG, PPM ou PPB. Valores: M\_M,V\_V,M\_V,MOL\_L,MOL\_KG,PPM,PPB.
\item[\texttt{notacao\_original\_fabricante}] string; N. Rótulo original, até 50 caracteres; não substitui limites numéricos.
\end{description}

\paragraph{Especificacao\_Reagente}\mbox{}\par
Caminho: \texttt{Resumo\_Reagente/\{resumoId\}/Especificacoes/\{id\}}.
\begin{description}
\item[\texttt{id}] docId string; O. Chave do documento, gerada pelo servidor; não duplicar como contador relacional.
\item[\texttt{id\_resumo\_reagente}] string; O. Identificador da entidade resumo reagente, cuja existência e escopo são verificados no servidor.
\item[\texttt{descricao}] string; O. Valor de descricao, validado pelo formulário e pelo servidor conforme regras do domínio. Máximo 150 caracteres.
\item[\texttt{fabricante}] string; N. Nome do fabricante da especificação comercial; não é o fornecedor do lote. Máximo 150 caracteres.
\item[\texttt{codigo\_produto\_fabricante}] string; N. Código comercial do fabricante para distinguir versões de produto sob o mesmo resumo. Máximo 100 caracteres.
\item[\texttt{grau\_pureza}] string; N. Grau declarado pelo fabricante, por exemplo P.A. ou HPLC; texto opcional, não percentual obrigatório. Máximo 30 caracteres.
\item[\texttt{densidade}] number; N. Densidade em g/ml; positiva e obrigatória para líquido, imutável após confirmação.
\item[\texttt{unidade\_de\_medida}] string enum; projeção física opcional derivada pelo servidor do estado do resumo, g/ml; não é fonte canônica nem campo editável.
\item[\texttt{classe\_inflamabilidade}] string enum; O. Classe selecionada para a especificação; não derivar apenas do nome comercial. Valores: NAO\_INFLAMAVEL,CLASSE\_1,CLASSE\_2,CLASSE\_3.
\item[\texttt{eh\_controlado\_pf}] boolean; O. Indicador cadastral de controle PF declarado pelo gestor; não é inferência automática nem parecer regulatório.
\item[\texttt{eh\_controlado\_eb}] boolean; O. Indicador cadastral de controle EB declarado pelo gestor, independente do indicador PF.
\item[\texttt{link\_fds\_fispq}] string; N. Endereço HTTPS opcional para a ficha de segurança; abrir com proteção de navegação externa. Máximo 255 caracteres.
\end{description}

\paragraph{Lote}\mbox{}\par
Caminho: \texttt{Lote/\{id\}}.
\begin{description}
\item[\texttt{id}] docId string; O. Chave do documento, gerada pelo servidor; não duplicar como contador relacional.
\item[\texttt{id\_resumo\_reagente}] string; O. Projeção física imutável do documento pai da especificação, preenchida no cadastro para permitir agregações por resumo (desvio de 3FN focado em leitura). A base de dados requer backfill.
\item[\texttt{id\_especificacao\_reagente}] string; O. Identificador da entidade especificacao reagente, cuja existência e escopo são verificados no servidor.
\item[\texttt{data\_aquisicao}] Timestamp; O. Data da aquisição do lote; não pode estar no futuro no cadastro de entrada já realizada.
\item[\texttt{qtd\_frascos\_comprados}] number; O. Quantidade adquirida, distinta da contagem materializada de frascos cadastrados.
\item[\texttt{nome\_fornecedor}] string; O. Fornecedor da aquisição; participa da chave composta do lote. Máximo 80 caracteres.
\item[\texttt{numero\_lote}] string; O. Identificação do lote do fornecedor/fabricante, distinta do código LCQUI de cada frasco. Máximo 100 caracteres.
\item[\texttt{nota\_fiscal}] string; O. Referência textual do documento de aquisição; não contém o binário da nota. Máximo 100 caracteres.
\item[\texttt{data\_fabricacao}] string YYYY-MM-DD; O. Data civil declarada de fabricação; não posterior à validade.
\item[\texttt{data\_validade}] string YYYY-MM-DD; O. Data civil de validade do lote; preenche sugestão de validade fechada do frasco.
\end{description}

\paragraph{Frasco\_Reagente}\mbox{}\par
Caminho: \texttt{Frasco\_Reagente/\{id\}}.
\begin{description}
\item[\texttt{id}] docId string; O. Chave do documento, gerada pelo servidor; não duplicar como contador relacional.
\item[\texttt{id\_almoxarifado}] string; O. Identificador da entidade almoxarifado, cuja existência e escopo são verificados no servidor.
\item[\texttt{id\_lote}] string; N. Identificador da entidade lote, cuja existência e escopo são verificados no servidor.
\item[\texttt{id\_especificacao\_reagente}] string; O. Identificador da entidade especificacao reagente, cuja existência e escopo são verificados no servidor. Projeção imutável obrigatória no Firestore.
\item[\texttt{id\_resumo\_reagente}] string; O. Identificador da entidade resumo reagente, cuja existência e escopo são verificados no servidor. Projeção imutável obrigatória no Firestore.
\item[\texttt{detalhe\_local\_armazenamento}] string; N. Valor de detalhe local armazenamento, validado pelo formulário e pelo servidor conforme regras do domínio.
\item[\texttt{codigo\_frasco}] string; O. Identificador LCQUI-N único, gerado pelo servidor. Unicidade garantida por chave determinística/registro de unicidade no servidor.
\item[\texttt{data\_ultima\_pesagem}] Timestamp; N. Data/instante de data ultima pesagem, com ordem temporal validada e autoria do relógio do servidor para eventos.
\item[\texttt{conteudo\_nominal}] number; N. Conteúdo inicial declarado em ml ou g conforme especificação. Se NULL, valor original do rótulo desconhecido ou ilegível (não usar 0 para isso).
\item[\texttt{peso\_no\_cadastrado}] number; O. Peso bruto inicial em g; imutável.
\item[\texttt{peso\_atual}] number; O. Último peso bruto confirmado em g.
\item[\texttt{peso\_frasco\_vazio}] number; N. Tara em g, calculada ou informada conforme modalidade.
\item[\texttt{medida\_usada}] number; O. Acumulado gravimétrico em g; não misturar com volume.
\item[\texttt{data\_abertura}] string YYYY-MM-DD; N. Data efetiva da abertura; não inventar data histórica desconhecida (Q05).
\item[\texttt{validade\_fechado}] string YYYY-MM-DD; N. Validade declarada antes da abertura; pode ser sugerida pelo lote, com origem preservada.
\item[\texttt{validade\_efetiva}] string YYYY-MM-DD; N. Prazo aplicável calculado com abertura e validade fechada; null quando não determinável.
\item[\texttt{validade\_desconhecida}] boolean; O. Ausência explícita de informação de validade; não significa validade ilimitada.
\item[\texttt{validade\_apos\_aberto\_dias}] number; N. Prazo inteiro positivo após abertura quando declarado; compõe o mínimo da validade efetiva.
\item[\texttt{estado\_fisico\_frasco}] string enum; O. Condição do recipiente/conteúdo; quarentena e vencimento são dimensões separadas. Valores: FECHADO,ABERTO,VAZIO,QUEBRADO,DESCARTADO,EXTRAVIADO.
\item[\texttt{disponibilidade}] string enum; O. Disponível ou emprestado; exige coerência com empréstimo ativo. Valores: DISPONIVEL,EMPRESTADO.
\item[\texttt{vencido}] boolean; O. Projeção de validade comparada ao relógio do servidor; reavaliar em cada operação crítica.
\item[\texttt{em\_quarentena}] boolean; O. Bloqueio operacional independente de abertura e vencimento; true exige motivo.
\item[\texttt{uso\_vencido\_autorizado}] boolean; O. Decisão expressa de gestor para uso excepcional conforme Q04, auditada; pesquisa exige TCR e justificativa.
\item[\texttt{detalhe\_status}] string; N. Justificativa obrigatória para quarentena e decisões que a exigem.
\item[\texttt{cadastrado\_em}] Timestamp; O. Data/instante de cadastrado em, com ordem temporal validada e autoria do relógio do servidor para eventos.
\item[\texttt{cadastrado\_por}] string; O. Valor de cadastrado por, validado pelo formulário e pelo servidor conforme regras do domínio.
\item[\texttt{eh\_higroscopico}] boolean; O. Snapshot físico denormalizado, imutável, preenchido pelo backend no cadastro a partir de Resumo\_Reagente.eh\_higroscopico (fonte canônica 3FN). Alterações posteriores do resumo não reescrevem frascos existentes. Não aceitar valor do cliente nem atualização direta. Permite aplicar Q06 na devolução sem leitura adicional do resumo. Frascos legados sem snapshot exigem reconciliação explícita antes da operação.
\item[\texttt{abertura\_historica\_desconhecida}] boolean; O, padrão false. True exige data\_abertura null e estado coerente com abertura anterior, sem inventar data.
\end{description}

\paragraph{Historico\_Frasco\_Reagente}\mbox{}\par
Caminho: \texttt{Historico\_Frasco\_Reagente/\{id\}}.
\begin{description}
\item[\texttt{id}] docId string; O. Chave do documento, gerada pelo servidor; não duplicar como contador relacional.
\item[\texttt{id\_frasco\_reagente}] string; O. Identificador da entidade frasco reagente, cuja existência e escopo são verificados no servidor.
\item[\texttt{id\_almoxarifado}] string; O. Identificador da entidade almoxarifado, cuja existência e escopo são verificados no servidor.
\item[\texttt{id\_gestor}] string; O. Identificador da entidade gestor, cuja existência e escopo são verificados no servidor.
\item[\texttt{tipo}] string enum; O. Valor de tipo, validado pelo formulário e pelo servidor conforme regras do domínio.
\item[\texttt{campo\_ajustado}] string; N. Nome permitido do atributo alterado no evento; não é caminho arbitrário para atualização de documento. Máximo 30 caracteres.
\item[\texttt{id\_emprestimo\_reagente}] string; N. Identificador da entidade emprestimo reagente, cuja existência e escopo são verificados no servidor.
\item[\texttt{peso\_anterior}] number; N. Peso bruto em g imediatamente antes do evento de ajuste/pesagem, quando aplicável.
\item[\texttt{peso\_novo}] number; N. Peso bruto em g confirmado no evento, quando aplicável.
\item[\texttt{medida\_ajustada}] number; N. Diferença operacional na unidade indicada, com sinal definido pelo tipo de evento; não somar automaticamente como consumo.
\item[\texttt{unidade\_medida\_ajustada}] string enum; N. Valor de unidade medida ajustada, validado pelo formulário e pelo servidor conforme regras do domínio. Valores: ml,g. (NULL se evento não for quantitativo).
\item[\texttt{timestamp}] Timestamp; O. Data/instante de timestamp, com ordem temporal validada e autoria do relógio do servidor para eventos.
\end{description}

\paragraph{Emprestimo\_Reagente}\mbox{}\par
Caminho: \texttt{Emprestimo\_Reagente/\{id\}}.
\begin{description}
\item[\texttt{id}] docId string; O. Chave do documento, gerada pelo servidor; não duplicar como contador relacional.
\item[\texttt{id\_frasco\_reagente}] string; O. Identificador da entidade frasco reagente, cuja existência e escopo são verificados no servidor.
\item[\texttt{id\_usuario\_retirou}] string; O. Identificador da entidade usuario retirou, cuja existência e escopo são verificados no servidor.
\item[\texttt{status}] string enum; O. Estado controlado pelo domínio; transições são validadas no servidor. Valores: EM\_USO,ATRASADO,DEVOLVIDO,DEVOLVIDO\_COM\_ATRASO,ENCERRADO\_EXTRAORDINARIO.
\item[\texttt{data\_retirada}] Timestamp; O. Data/instante de data retirada, com ordem temporal validada e autoria do relógio do servidor para eventos.
\item[\texttt{data\_devolucao\_prevista}] string YYYY-MM-DD; O. Data civil cujo prazo termina no fim do dia no fuso institucional.
\item[\texttt{data\_devolucao\_efetuada}] Timestamp; N. Data/instante de data devolucao efetuada, com ordem temporal validada e autoria do relógio do servidor para eventos.
\item[\texttt{id\_local\_usado}] string; O. Identificador da entidade local usado, cuja existência e escopo são verificados no servidor.
\item[\texttt{id\_usuario\_devolveu}] string; N. Identificador da entidade usuario devolveu, cuja existência e escopo são verificados no servidor.
\item[\texttt{peso\_saida}] number; O. Peso bruto em g medido na retirada; referência imutável para o consumo deste empréstimo.
\item[\texttt{peso\_retorno}] number; N. Peso bruto em g medido na devolução; null antes do retorno e validado contra tara/saída.
\item[\texttt{medida\_utilizada}] number; N. Consumo não negativo calculado na unidade do empréstimo; massa = max(0, saída menos retorno), conversão em mL só para líquido com densidade histórica positiva.
\item[\texttt{unidade\_medida\_utilizada}] string enum; O. Valor de unidade medida utilizada, validado pelo formulário e pelo servidor conforme regras do domínio. Valores: ml,g.
\item[\texttt{peso\_perda\_evaporacao}] number; O. Perda em g registrada separadamente de consumo; não inferir arbitrariamente.
\item[\texttt{uso\_vencido\_aceito}] boolean; O. Ciência explícita do empréstimo específico; não concede nem substitui autorização do gestor.
\item[\texttt{finalidade\_uso}] string enum; O. Finalidade validada na retirada. Valores: AULA\_PRATICA, DEMONSTRACAO, PESQUISA\_TCC\_POS, ESTUDO\_DEGRADACAO\_RESIDUOS. Uso excepcional de pesquisa exige rastreabilidade TCR de Q04.
\item[\texttt{auto\_atendimento}] boolean; O, padrão false. Calculado pelo backend conforme Q14, justificativa e notificação à chefia obrigatórias quando true.
\item[\texttt{justificativa\_metodologica}] string; C no uso excepcional Q04, 20--2000 caracteres não vazios.
\item[\texttt{tcr\_versao}] string; C em Q04, até 100. Versão institucional aceita, validada pelo backend.
\item[\texttt{tcr\_aceito\_em}] Timestamp; C em Q04. Instante servidor do aceite.
\item[\texttt{tcr\_aceito\_por}] string; C em Q04. UID autenticado do solicitante/retirante, não do gestor que atende.
\item[\texttt{tcr\_auditoria\_id}] string; C em Q04. Evento imutável em Registro\_de\_Auditoria com frasco, operação, finalidade, versão, instante e aceitante; backend verifica correspondência e consumo único do aceite.
\end{description}

\paragraph{Turma}\mbox{}\par
Caminho: \texttt{Turma/\{id\}}.
\begin{description}
\item[\texttt{id}] docId string; O. Chave do documento, gerada pelo servidor; não duplicar como contador relacional.
\item[\texttt{id\_materia}] string; O. Identificador da entidade materia, cuja existência e escopo são verificados no servidor.
\item[\texttt{id\_professor}] string; O. Identificador da entidade professor, cuja existência e escopo são verificados no servidor.
\item[\texttt{status}] string enum; O. Estado controlado pelo domínio; transições são validadas no servidor. Valores: Ativo, Arquivada.
\item[\texttt{nome\_turma}] string; O. Nome de apresentação da turma, sem inferir ano/semestre a partir do texto. Máximo 100 caracteres.
\item[\texttt{ano}] number; O. Ano letivo da turma ou ano da competência do agregado mensal, conforme documento.
\item[\texttt{semestre}] number; O. Semestre letivo inteiro, exclusivamente 1 ou 2.
\item[\texttt{capacidade}] number; O. Quantidade máxima ordinária de alunos; inteiro positivo, validado na transação de ingresso.
\item[\texttt{codigo\_turma}] string; O. Código único de ingresso, não concede acesso sem validação. Máximo 20 caracteres. Unicidade garantida por chave determinística/registro de unicidade no servidor.
\item[\texttt{data\_criacao}] Timestamp; O. Data/instante de data criacao, com ordem temporal validada e autoria do relógio do servidor para eventos.
\end{description}

\paragraph{Historico\_Alunos\_Turma}\mbox{}\par
Caminho: \texttt{Turma/\{id\}/HistoricoAlunos/\{eventoId\}}.
\begin{description}
\item[\texttt{id}] docId string; O. Chave do documento, gerada pelo servidor; não duplicar como contador relacional.
\item[\texttt{id\_turma}] string; O. Identificador da entidade turma, cuja existência e escopo são verificados no servidor.
\item[\texttt{tipo}] string enum; O. Valor de tipo, validado pelo formulário e pelo servidor conforme regras do domínio. Valores: inclusao\_aluno, exclusao\_aluno.
\item[\texttt{id\_aluno}] string; O. Identificador da entidade aluno, cuja existência e escopo são verificados no servidor.
\item[\texttt{timestamp}] Timestamp; O. Data/instante de timestamp, com ordem temporal validada e autoria do relógio do servidor para eventos.
\end{description}

\paragraph{Historico\_Posts\_Turma}\mbox{}\par
Caminho: \texttt{Turma/\{id\}/Posts/\{postId\}/Historico/\{eventoId\}}.
\begin{description}
\item[\texttt{id}] docId string; O. Chave do documento, gerada pelo servidor; não duplicar como contador relacional.
\item[\texttt{id\_post}] string; O. Identificador da entidade post, cuja existência e escopo são verificados no servidor.
\item[\texttt{novo\_titulo}] string; N. Valor novo titulo, restrito ao campo correspondente da entidade de origem; preservar ausência legítima. Máximo 150 caracteres.
\item[\texttt{novo\_id\_roteiro}] string; N. Valor novo id roteiro, restrito ao campo correspondente da entidade de origem; preservar ausência legítima.
\item[\texttt{nova\_descricao}] string; N. Valor nova descricao, restrito ao campo correspondente da entidade de origem; preservar ausência legítima.
\item[\texttt{antigo\_titulo}] string; N. Valor antigo titulo, restrito ao campo correspondente da entidade de origem; preservar ausência legítima. Máximo 150 caracteres.
\item[\texttt{antigo\_id\_roteiro}] string; N. Valor antigo id roteiro, restrito ao campo correspondente da entidade de origem; preservar ausência legítima.
\item[\texttt{antiga\_descricao}] string; N. Valor antiga descricao, restrito ao campo correspondente da entidade de origem; preservar ausência legítima.
\item[\texttt{timestamp}] Timestamp; O. Data/instante de timestamp, com ordem temporal validada e autoria do relógio do servidor para eventos.
\end{description}

\paragraph{Historico\_Comentario}\mbox{}\par
Caminho: \texttt{Turma/\{id\}/Posts/\{postId\}/Comentarios/\{id\}/Historico/\{eventoId\}}.
\begin{description}
\item[\texttt{id}] docId string; O. Chave do documento, gerada pelo servidor; não duplicar como contador relacional.
\item[\texttt{id\_comentario}] string; O. Identificador da entidade comentario, cuja existência e escopo são verificados no servidor.
\item[\texttt{novo\_texto}] string; O. Valor novo texto, restrito ao campo correspondente da entidade de origem; preservar ausência legítima.
\item[\texttt{texto\_antigo}] string; O. Valor de texto antigo, validado pelo formulário e pelo servidor conforme regras do domínio.
\item[\texttt{editado\_em}] Timestamp; O. Data/instante de editado em, com ordem temporal validada e autoria do relógio do servidor para eventos.
\item[\texttt{editado\_por}] string; O. Valor de editado por, validado pelo formulário e pelo servidor conforme regras do domínio.
\end{description}

\paragraph{Roteiro\_Experimento}\mbox{}\par
Caminho: \texttt{Roteiro\_Experimento/\{id\}}.
\begin{description}
\item[\texttt{id}] docId string; O. Chave do documento, gerada pelo servidor; não duplicar como contador relacional.
\item[\texttt{id\_professor\_upload}] string; O. Identificador da entidade professor upload, cuja existência e escopo são verificados no servidor.
\item[\texttt{file\_url}] string; O. Referência do PDF; download exige revalidar acesso.
\item[\texttt{nome}] string; O. Nome de apresentação; não substitui a chave do documento. Máximo 150 caracteres.
\item[\texttt{descricao}] string; O. Valor de descricao, validado pelo formulário e pelo servidor conforme regras do domínio.
\item[\texttt{professores\_compartilhados}] array de strings; O, padrão vazio. UIDs autorizados, geridos exclusivamente pelo servidor com arrayUnion/arrayRemove; sem escrita direta cliente.
\end{description}

A associação Roteiro\_Professor\_Compartilhado existe somente no modelo relacional 3FN; no Firestore é o array ACL professores\_compartilhados em Roteiro\_Experimento. Consultar Meus por id\_professor\_upload == uid; Compartilhados comigo por array-contains no UID; Eu compartilhei filtra os roteiros próprios carregados com ACL não vazia.

\paragraph{Professor}\mbox{}\par
Caminho: \texttt{Professor/\{id\}}.
\begin{description}
\item[\texttt{id\_usuario}] string; O. Identificador da entidade usuario, cuja existência e escopo são verificados no servidor.
\item[\texttt{centro}] string enum; O. Centro institucional do professor: CCT, CCTA, CBB ou CCH. Valores: CCT, CCTA, CBB, CCH.
\item[\texttt{laboratorio}] string; O. Laboratório de lotação do professor; não fornece automaticamente prédio/andar/sala. Máximo 100 caracteres.
\end{description}

\paragraph{Materia}\mbox{}\par
Caminho: \texttt{Materia/\{id\}}.
\begin{description}
\item[\texttt{id}] docId string; O. Chave do documento, gerada pelo servidor; não duplicar como contador relacional.
\item[\texttt{nome}] string; O. Nome de apresentação; não substitui a chave do documento. Máximo 100 caracteres.
\item[\texttt{codigo\_materia}] string; O. Código acadêmico único da matéria, normalizado de modo uniforme. Máximo 10 caracteres.
\end{description}

\paragraph{Professor\_x\_Materia}\mbox{}\par
Caminho: \texttt{Professor\_x\_Materia/\{id\}}.
\begin{description}
\item[\texttt{id\_professor}] string; O. Identificador da entidade professor, cuja existência e escopo são verificados no servidor.
\item[\texttt{id\_materia}] string; O. Identificador da entidade materia, cuja existência e escopo são verificados no servidor.
\end{description}

\paragraph{Aluno}\mbox{}\par
Caminho: \texttt{Aluno/\{id\}}.
\begin{description}
\item[\texttt{id\_usuario}] string; O. Identificador da entidade usuario, cuja existência e escopo são verificados no servidor.
\item[\texttt{numero\_matricula}] string; O. Matrícula institucional textual, com zeros iniciais; obrigatória e única em Aluno, opcional e não única em Convite. Máximo 20 caracteres.
\end{description}

\paragraph{Bolsista}\mbox{}\par
Caminho: \texttt{Bolsista/\{id\}}.
\begin{description}
\item[\texttt{id\_usuario}] string; O. Identificador da entidade usuario, cuja existência e escopo são verificados no servidor.
\end{description}

\paragraph{Convite\_Aluno}\mbox{}\par
Caminho: \texttt{Convite\_Aluno/\{id\}}.
\begin{description}
\item[\texttt{id}] docId string; O. Chave do documento, gerada pelo servidor; não duplicar como contador relacional.
\item[\texttt{id\_turma}] string; N. Identificador da entidade turma, cuja existência e escopo são verificados no servidor.
\item[\texttt{email}] string; O. Endereço normalizado para identidade/convite; não divulgar em catálogo público interno. Máximo 150 caracteres.
\item[\texttt{convidado\_em}] Timestamp; O. Data/instante de convidado em, com ordem temporal validada e autoria do relógio do servidor para eventos.
\item[\texttt{status}] string enum; O. Estado controlado pelo domínio; transições são validadas no servidor. Valores: pendente, aceitado, expirado.
\item[\texttt{expira\_em}] Timestamp; O. Instante de expiração de convite/notificação; não equivale a autorização de excluir auditoria.
\item[\texttt{convidado\_por}] string; O. Valor de convidado por, validado pelo formulário e pelo servidor conforme regras do domínio.
\item[\texttt{numero\_matricula}] string; N. Matrícula institucional textual, com zeros iniciais; obrigatória e única em Aluno, opcional e não única em Convite. Máximo 20 caracteres.
\item[\texttt{token\_hash}] string; O. SHA-256 do token aleatório de 32 bytes (CSPRNG); nunca token em claro. Reenvio substitui hash e expiração atomicamente; comparação constante e aceitação única no servidor.
\item[\texttt{ultimo\_reenvio\_por}] string; N. UID servidor do operador do último reenvio. DocId derivado por HMAC-SHA-256 de contexto e e-mail normalizado com segredo servidor; não usar hash simples público de e-mail.
\end{description}

\paragraph{Notificacao}\mbox{}\par
Caminho: \texttt{Usuarios/\{uid\}/Notificacoes/\{id\}}.
\begin{description}
\item[\texttt{id}] docId string; O. Chave do documento, gerada pelo servidor; não duplicar como contador relacional.
\item[\texttt{id\_destinatario}] string; O. Identificador da entidade destinatario, cuja existência e escopo são verificados no servidor.
\item[\texttt{papel\_destinatario}] string enum; O. Contexto de exibição da notificação, sem conceder esse papel ao usuário. Inclui Aluno, Professor, Bolsista, Gestor\_Bens\_Patrimoniais e Gestor\_Almoxarifado.
\item[\texttt{tipo}] string enum; O. Valor de tipo, validado pelo formulário e pelo servidor conforme regras do domínio.
\item[\texttt{id\_quem\_fez\_acao}] string; N. Identificador da entidade quem fez acao, cuja existência e escopo são verificados no servidor.
\item[\texttt{id\_turma}] string; N. Identificador da entidade turma, cuja existência e escopo são verificados no servidor.
\item[\texttt{quantidade}] number; N. Quantidade agregada de ocorrências agrupadas na notificação, quando aplicável.
\item[\texttt{entidade\_alvo}] string; O. Tipo autorizado de recurso para navegação; mapear a rotas conhecidas, não URL arbitrária. Máximo 40 caracteres.
\item[\texttt{id\_alvo}] string; O. ID do recurso da notificação; acesso é revalidado ao abrir. Máximo 200 caracteres.
\item[\texttt{lida}] boolean; O. Estado de leitura da notificação; true preserva o registro em vez de excluí-lo.
\item[\texttt{lida\_em}] Timestamp; N. Instante servidor de marcação de leitura; null enquanto não lida.
\item[\texttt{emitida\_em}] Timestamp; O. Data/instante de emitida em, com ordem temporal validada e autoria do relógio do servidor para eventos.
\item[\texttt{expira\_em}] Timestamp; O. Instante de expiração de convite/notificação; não equivale a autorização de excluir auditoria.
\end{description}

\paragraph{Aluno\_x\_Turma}\mbox{}\par
Caminho: \texttt{Turma/\{id\}/Alunos/\{uid\}}.
\begin{description}
\item[\texttt{id\_aluno}] string; O. Identificador da entidade aluno, cuja existência e escopo são verificados no servidor.
\item[\texttt{id\_turma}] string; O. Identificador da entidade turma, cuja existência e escopo são verificados no servidor.
\end{description}

\paragraph{Local}\mbox{}\par
Caminho: \texttt{Local/\{id\}}.
\begin{description}
\item[\texttt{id}] docId string; O. Chave do documento, gerada pelo servidor; não duplicar como contador relacional.
\item[\texttt{predio}] string; O. Identificador textual do prédio no cadastro de locais. Máximo 30 caracteres.
\item[\texttt{andar}] string; O. Identificador textual do piso; permite térreo/subsolo, sem impor número inteiro. Máximo 10 caracteres.
\item[\texttt{sala}] string; O. Identificador textual da sala, único no contexto prédio/andar. Máximo 30 caracteres.
\item[\texttt{criado\_por}] string; O. Valor de criado por, validado pelo formulário e pelo servidor conforme regras do domínio.
\item[\texttt{criado\_em}] Timestamp; O. Data/instante de criado em, com ordem temporal validada e autoria do relógio do servidor para eventos.
\end{description}

\paragraph{Post}\mbox{}\par
Caminho: \texttt{Turma/\{id\}/Posts/\{postId\}}.
\begin{description}
\item[\texttt{id}] docId string; O. Chave do documento, gerada pelo servidor; não duplicar como contador relacional.
\item[\texttt{id\_professor}] string; O. Identificador da entidade professor, cuja existência e escopo são verificados no servidor.
\item[\texttt{id\_turma}] string; O. Identificador da entidade turma, cuja existência e escopo são verificados no servidor.
\item[\texttt{id\_roteiro\_experimento}] string; N. Identificador da entidade roteiro experimento, cuja existência e escopo são verificados no servidor.
\item[\texttt{titulo}] string; O. Título obrigatório de publicação, exibido no feed e na navegação de notificações. Máximo 150 caracteres.
\item[\texttt{descricao}] string; O. Valor de descricao, validado pelo formulário e pelo servidor conforme regras do domínio.
\item[\texttt{criado\_em}] Timestamp; O. Data/instante de criado em, com ordem temporal validada e autoria do relógio do servidor para eventos.
\item[\texttt{roteiro\_anexo}] map; N. Snapshot imutável da publicação: id\_roteiro (string O), nome\_arquivo (string O, até 150), tamanho\_bytes (inteiro positivo O), storage\_path (string O). Backend valida acesso ao roteiro antes de publicar. Revogação posterior da ACL não remove o anexo histórico. Aluno recebe os metadados no feed; download via obterUrlDownloadRoteiro valida vínculo atual e emite URL assinada por 15 minutos, sem leitura cliente da coleção raiz.
\end{description}

\paragraph{Comentario}\mbox{}\par
Caminho: \texttt{Turma/\{id\}/Posts/\{postId\}/Comentarios/\{id\}}.
\begin{description}
\item[\texttt{id}] docId string; O. Chave do documento, gerada pelo servidor; não duplicar como contador relacional.
\item[\texttt{id\_post}] string; O. Identificador da entidade post, cuja existência e escopo são verificados no servidor.
\item[\texttt{id\_usuario}] string; O. Identificador da entidade usuario, cuja existência e escopo são verificados no servidor.
\item[\texttt{texto}] string; O. Conteúdo textual do comentário, escapado na renderização; rejeitar texto vazio após trim.
\item[\texttt{criado\_em}] Timestamp; O. Data/instante de criado em, com ordem temporal validada e autoria do relógio do servidor para eventos.
\item[\texttt{moderado}] boolean; O, padrão false. Visibilidade conforme DP-C02.
\item[\texttt{motivo\_moderacao}] string; C quando moderado, até 2000 caracteres não vazios.
\item[\texttt{moderado\_por}] string; C quando moderado. UID autenticado do moderador, preenchido pelo servidor.
\end{description}

\paragraph{Registro\_de\_Auditoria}\mbox{}\par
Caminho: \texttt{Registro\_de\_Auditoria/\{id\}}.
\begin{description}
\item[\texttt{id}] docId string; O. Chave do documento, gerada pelo servidor; não duplicar como contador relacional.
\item[\texttt{id\_usuario}] string; O. Identificador da entidade usuario, cuja existência e escopo são verificados no servidor.
\item[\texttt{acao}] string; O. Nome da ação auditada; distinguir solicitação, conclusão e rejeição quando relevante. Máximo 200 caracteres.
\item[\texttt{tipo\_entidade\_sofre\_acao}] string enum; O. Valor de tipo entidade sofre acao, validado pelo formulário e pelo servidor conforme regras do domínio.
\item[\texttt{id\_do\_objeto\_da\_entidade}] string; O. Identificador da entidade do objeto da entidade, cuja existência e escopo são verificados no servidor. Máximo 200 caracteres.
\item[\texttt{acao\_feita\_em}] Timestamp; O. Data/instante de acao feita em, com ordem temporal validada e autoria do relógio do servidor para eventos.
\item[\texttt{metadata}] map; O. Mapa auditável e limitado ao evento, sem senha, token ou conteúdo binário.
\end{description}

\paragraph{Lote\_Materializado}\mbox{}\par
Caminho: \texttt{Lote\_Materializado/\{id\}}.
\begin{description}
\item[\texttt{id}] docId string; O. Chave do documento, gerada pelo servidor; não duplicar como contador relacional.
\item[\texttt{id\_lote}] string; O. Identificador da entidade lote, cuja existência e escopo são verificados no servidor. Unicidade garantida por chave determinística/registro de unicidade no servidor.
\item[\texttt{qtd\_frascos\_cadastrados}] number; O. Contagem de frascos cadastrados, inteira não negativa; não editável pelo cliente quando materializada.
\item[\texttt{ultimo\_reconciliador}] Timestamp; O (se já reconciliado). Marca temporal absoluta, preenchida/atualizada pelo job de reconciliação. Eventos de trigger gerados antes deste instante são ignorados, evitando que eventos assíncronos retardados somem duplamente sobre a contagem absoluta já aferida.
\end{description}

\paragraph{Auditoria\_Reconciliacao}\mbox{}\par
Caminho: \texttt{Auditoria\_Reconciliacao/\{id\}}.
\begin{description}
\item[\texttt{id}] docId string; O. Chave gerada pelo servidor para o registro da cura.
\item[\texttt{id\_lote}] string; O. Lote alvo da reconciliação onde ocorreu divergência.
\item[\texttt{qtd\_anterior}] number; O. Contagem que constava no documento \texttt{Lote\_Materializado} antes da correção (drift).
\item[\texttt{qtd\_corrigida}] number; O. Contagem absoluta efetivamente encontrada no banco de dados e aplicada (count real).
\item[\texttt{data}] Timestamp; O. Instante exato em que a correção transacional foi aplicada e a marca d'água elevada.
\end{description}

\paragraph{Resumo\_Almoxarifado\_Diario}\mbox{}\par
Caminho: \texttt{Resumo\_Almoxarifado\_Diario/\{id\}}.
\begin{description}
\item[\texttt{id}] docId string; O. Chave do documento, gerada pelo servidor; não duplicar como contador relacional.
\item[\texttt{data}] string YYYY-MM-DD; O. Data/instante de data, com ordem temporal validada e autoria do relógio do servidor para eventos.
\item[\texttt{id\_almoxarifado}] string; O. Identificador da entidade almoxarifado, cuja existência e escopo são verificados no servidor.
\item[\texttt{qtd\_frascos\_fechados\_no\_fim\_do\_dia}] number; O. Contagem de frascos fechados no fim do dia, inteira não negativa; não editável pelo cliente quando materializada.
\item[\texttt{qtd\_frascos\_abertos\_no\_fim\_do\_dia}] number; O. Contagem de frascos abertos no fim do dia, inteira não negativa; não editável pelo cliente quando materializada.
\item[\texttt{qtd\_frascos\_emprestados\_no\_fim\_do\_dia}] number; O. Contagem de frascos emprestados no fim do dia, inteira não negativa; não editável pelo cliente quando materializada.
\item[\texttt{qtd\_frascos\_cadastrados\_fechados\_durante\_o\_dia}] number; O. Contagem de frascos cadastrados fechados durante o dia, inteira não negativa; não editável pelo cliente quando materializada.
\item[\texttt{qtd\_frascos\_cadastrados\_abertos\_durante\_o\_dia}] number; O. Contagem de frascos cadastrados abertos durante o dia, inteira não negativa; não editável pelo cliente quando materializada.
\item[\texttt{qtd\_frascos\_emprestados\_durante\_o\_dia}] number; O. Contagem de frascos emprestados durante o dia, inteira não negativa; não editável pelo cliente quando materializada.
\item[\texttt{qtd\_frascos\_devolvidos\_durante\_o\_dia}] number; O. Contagem de frascos devolvidos durante o dia, inteira não negativa; não editável pelo cliente quando materializada.
\item[\texttt{volume\_total\_usado\_nos\_frascos\_devolvidos\_durante\_o\_dia}] number; O. Métrica de volume total usado nos frascos devolvidos durante o dia, não negativa e calculada a partir dos eventos/saldos do domínio.
\item[\texttt{volume\_total\_disponivel\_ml}] number; O. Métrica de volume total disponivel ml, não negativa e calculada a partir dos eventos/saldos do domínio.
\item[\texttt{volume\_utilizado\_no\_dia\_ml}] number; O. Métrica de volume utilizado no dia ml, não negativa e calculada a partir dos eventos/saldos do domínio.
\item[\texttt{massa\_total\_disponivel\_g}] number; O. Métrica de massa total disponivel g, não negativa e calculada a partir dos eventos/saldos do domínio.
\item[\texttt{massa\_utilizada\_no\_dia\_g}] number; O. Métrica de massa utilizada no dia g, não negativa e calculada a partir dos eventos/saldos do domínio.
\item[\texttt{created\_at}] Timestamp; O. Instante servidor da primeira criação da materialização.
\item[\texttt{updated\_at}] Timestamp; O. Instante servidor da última atualização da materialização; informa atualidade do resultado.
\end{description}

\paragraph{Resumo\_Reagente\_Diario}\mbox{}\par
Caminho: \texttt{Resumo\_Reagente\_Diario/\{id\}}.
\begin{description}
\item[\texttt{id}] docId string; O. Chave do documento, gerada pelo servidor; não duplicar como contador relacional.
\item[\texttt{data}] string YYYY-MM-DD; O. Data/instante de data, com ordem temporal validada e autoria do relógio do servidor para eventos.
\item[\texttt{id\_resumo\_reagente}] string; O. Identificador da entidade resumo reagente, cuja existência e escopo são verificados no servidor.
\item[\texttt{id\_almoxarifado}] string; O. Identificador da entidade almoxarifado, cuja existência e escopo são verificados no servidor.
\item[\texttt{qtd\_frascos\_fechados\_disponiveis}] number; O. Contagem de frascos fechados disponiveis, inteira não negativa; não editável pelo cliente quando materializada.
\item[\texttt{qtd\_frascos\_abertos\_disponiveis}] number; O. Contagem de frascos abertos disponiveis, inteira não negativa; não editável pelo cliente quando materializada.
\item[\texttt{qtd\_frascos\_em\_uso}] number; O. Contagem de frascos em uso, inteira não negativa; não editável pelo cliente quando materializada.
\item[\texttt{qtd\_frascos\_vazios}] number; O. Contagem de frascos vazios, inteira não negativa; não editável pelo cliente quando materializada.
\item[\texttt{qtd\_frascos\_descartados\_dia}] number; O. Contagem de frascos descartados dia, inteira não negativa; não editável pelo cliente quando materializada.
\item[\texttt{volume\_total\_disponivel\_ml}] number; O. Métrica de volume total disponivel ml, não negativa e calculada a partir dos eventos/saldos do domínio.
\item[\texttt{volume\_utilizado\_no\_dia\_ml}] number; O. Métrica de volume utilizado no dia ml, não negativa e calculada a partir dos eventos/saldos do domínio.
\item[\texttt{volume\_evaporado\_no\_dia\_ml}] number; O. Métrica de volume evaporado no dia ml, não negativa e calculada a partir dos eventos/saldos do domínio.
\item[\texttt{massa\_total\_disponivel\_g}] number; O. Métrica de massa total disponivel g, não negativa e calculada a partir dos eventos/saldos do domínio.
\item[\texttt{massa\_utilizada\_no\_dia\_g}] number; O. Métrica de massa utilizada no dia g, não negativa e calculada a partir dos eventos/saldos do domínio.
\item[\texttt{massa\_evaporada\_no\_dia\_g}] number; O. Métrica de massa evaporada no dia g, não negativa e calculada a partir dos eventos/saldos do domínio.
\item[\texttt{created\_at}] Timestamp; O. Instante servidor da primeira criação da materialização.
\item[\texttt{updated\_at}] Timestamp; O. Instante servidor da última atualização da materialização; informa atualidade do resultado.
\end{description}

\paragraph{Resumo\_Bem\_Patrimonial\_Diario}\mbox{}\par
Caminho: \texttt{Resumo\_Bem\_Patrimonial\_Diario/\{id\}}.
\begin{description}
\item[\texttt{id}] docId string; O. Chave do documento, gerada pelo servidor; não duplicar como contador relacional.
\item[\texttt{data}] string YYYY-MM-DD; O. Data/instante de data, com ordem temporal validada e autoria do relógio do servidor para eventos.
\item[\texttt{id\_local}] string; N. Identificador da entidade local, cuja existência e escopo são verificados no servidor.
\item[\texttt{qtd\_bens\_ativos\_no\_fim\_do\_dia}] number; O. Contagem de bens ativos no fim do dia, inteira não negativa; não editável pelo cliente quando materializada.
\item[\texttt{qtd\_bens\_inservivel\_no\_fim\_do\_dia}] number; O. Contagem de bens inservivel no fim do dia, inteira não negativa; não editável pelo cliente quando materializada.
\item[\texttt{qtd\_bens\_dado\_baixa\_no\_fim\_do\_dia}] number; O. Contagem de bens dado baixa no fim do dia, inteira não negativa; não editável pelo cliente quando materializada.
\item[\texttt{qtd\_bens\_cadastrados\_durante\_o\_dia}] number; O. Contagem de bens cadastrados durante o dia, inteira não negativa; não editável pelo cliente quando materializada.
\item[\texttt{qtd\_bens\_marcados\_inservivel\_durante\_o\_dia}] number; O. Contagem de bens marcados inservivel durante o dia, inteira não negativa; não editável pelo cliente quando materializada.
\item[\texttt{qtd\_bens\_dados\_baixa\_durante\_o\_dia}] number; O. Contagem de bens dados baixa durante o dia, inteira não negativa; não editável pelo cliente quando materializada.
\item[\texttt{qtd\_requisicoes\_abertas\_durante\_o\_dia}] number; O. Contagem de requisicoes abertas durante o dia, inteira não negativa; não editável pelo cliente quando materializada.
\item[\texttt{qtd\_requisicoes\_aprovadas\_durante\_o\_dia}] number; O. Contagem de requisicoes aprovadas durante o dia, inteira não negativa; não editável pelo cliente quando materializada.
\item[\texttt{qtd\_requisicoes\_rejeitadas\_durante\_o\_dia}] number; O. Contagem de requisicoes rejeitadas durante o dia, inteira não negativa; não editável pelo cliente quando materializada.
\item[\texttt{created\_at}] Timestamp; O. Instante servidor da primeira criação da materialização.
\item[\texttt{updated\_at}] Timestamp; O. Instante servidor da última atualização da materialização; informa atualidade do resultado.
\end{description}

\paragraph{Atividade\_Gestor\_Almoxarifado\_Mensal}\mbox{}\par
Caminho: \texttt{Atividade\_Gestor\_Almoxarifado\_Mensal/\{id\}}.
\begin{description}
\item[\texttt{id}] docId string; O. Chave do documento, gerada pelo servidor; não duplicar como contador relacional.
\item[\texttt{id\_gestor}] string; O. Identificador da entidade gestor, cuja existência e escopo são verificados no servidor.
\item[\texttt{ano}] number; O. Ano letivo da turma ou ano da competência do agregado mensal, conforme documento.
\item[\texttt{mes}] number; O. Mês da competência mensal, inteiro de 1 a 12.
\item[\texttt{qtd\_reagentes\_adicionados}] number; O. Contagem de reagentes adicionados, inteira não negativa; não editável pelo cliente quando materializada.
\item[\texttt{qtd\_frascos\_adicionados}] number; O. Contagem de frascos adicionados, inteira não negativa; não editável pelo cliente quando materializada.
\item[\texttt{qtd\_movimentacoes\_registradas}] number; O. Contagem de movimentacoes registradas, inteira não negativa; não editável pelo cliente quando materializada.
\item[\texttt{created\_at}] Timestamp; O. Instante servidor da primeira criação da materialização.
\item[\texttt{updated\_at}] Timestamp; O. Instante servidor da última atualização da materialização; informa atualidade do resultado.
\end{description}

\paragraph{Atividade\_Gestor\_Bens\_Patrimoniais\_Mensal}\mbox{}\par
Caminho: \texttt{Atividade\_Gestor\_Bens\_Patrimoniais\_Mensal/\{id\}}.
\begin{description}
\item[\texttt{id}] docId string; O. Chave do documento, gerada pelo servidor; não duplicar como contador relacional.
\item[\texttt{id\_gestor}] string; O. Identificador da entidade gestor, cuja existência e escopo são verificados no servidor.
\item[\texttt{ano}] number; O. Ano letivo da turma ou ano da competência do agregado mensal, conforme documento.
\item[\texttt{mes}] number; O. Mês da competência mensal, inteiro de 1 a 12.
\item[\texttt{qtd\_bens\_adicionados}] number; O. Contagem de bens adicionados, inteira não negativa; não editável pelo cliente quando materializada.
\item[\texttt{qtd\_bens\_editados}] number; O. Contagem de bens editados, inteira não negativa; não editável pelo cliente quando materializada.
\item[\texttt{qtd\_requisicoes\_aprovadas}] number; O. Contagem de requisicoes aprovadas, inteira não negativa; não editável pelo cliente quando materializada.
\item[\texttt{qtd\_requisicoes\_rejeitadas}] number; O. Contagem de requisicoes rejeitadas, inteira não negativa; não editável pelo cliente quando materializada.
\item[\texttt{created\_at}] Timestamp; O. Instante servidor da primeira criação da materialização.
\item[\texttt{updated\_at}] Timestamp; O. Instante servidor da última atualização da materialização; informa atualidade do resultado.
\end{description}

\paragraph{Complementos físicos obrigatórios e fontes de verdade.}
\begin{description}
\item[Usuarios] \texttt{ativo} (boolean, O), \texttt{versao\_permissoes} (inteiro, O), \texttt{claims\_pendentes} (boolean, O), \texttt{atualizado\_em} (Timestamp, O): controle servidor de ativação e sincronização de autorização. Não replicar papéis editáveis no navegador.
\item[Perfis] docId = UID em Chefe\_Geral, Professor, Aluno, Bolsista e gestores; \texttt{id\_usuario} é esse UID. \texttt{nome}, \texttt{email} e \texttt{letra\_inicial}, se usados em busca de perfil, são projeções de Usuarios, com leitura restrita e atualização pelo servidor.
\item[Almoxarifado] \texttt{ativo} (boolean, O) e \texttt{qtd\_gestores\_ativos} (inteiro não negativo, O) são necessários à ativação e revogação concorrente. O contador é derivado dos vínculos; ativo exige contador positivo.
\item[Turma] \texttt{qtd\_alunos} (inteiro não negativo, O) é contador transacional do vínculo ativo, não count lido fora da transação. \texttt{versao} (inteiro positivo, O) permite detectar edições concorrentes.
\item[Espelho Aluno--Turma] \texttt{Usuarios/uid/Turmas/turmaId}: \texttt{id\_turma}, \texttt{id\_professor}, \texttt{id\_materia}, \texttt{nome\_turma}, \texttt{nome\_materia} (strings O), \texttt{ano}, \texttt{semestre} (inteiros O), \texttt{status} (enum O), \texttt{ingressou\_em} (Timestamp O). Todos derivados de Turma/Matéria ou evento de ingresso; não usar espelho desatualizado como autorização. Vínculo canônico em Turma/id/Alunos/uid contém \texttt{id\_aluno} (string O), \texttt{ingressou\_em} (Timestamp O); nome/foto para colegas são projeções mínimas, sem matrícula/e-mail.
\item[Reagentes] Resumo possui \texttt{letra\_inicial} (string O, derivada do nome). O estado físico escalar e a higroscopicidade são canônicos no resumo; filtro por igualdade. Especificação embute \texttt{composicao} (array de maps com os campos da composição; IDs de substâncias únicos). O resumo pai fornece seu ID; manter ID de resumo em Lote e Frasco para resolver o caminho completo da especificação.
\item[Lote e Frasco] \texttt{id\_resumo\_reagente}, \texttt{id\_especificacao\_reagente}, \texttt{nome\_reagente} (strings O) são projeções servidor; em Frasco, especificação permanece resolvida mesmo com lote, como desnormalização deliberada. \texttt{unidade\_medida} (ml/g, O), \texttt{localizacao\_detalhada} (string até 500, N), \texttt{versao} (inteiro O), \texttt{ultima\_pesagem\_em} (Timestamp, N) e \texttt{modalidade\_cadastro} (FECHADO/CONHECE\_TARA/ESTIMA\_VOLUME/ESTIMA\_MASSA, O) sustentam UI-06. \texttt{id\_emprestimo\_ativo} (string N) é ponteiro transacional, limpo no retorno. \texttt{id\_almoxarifado} é fonte da localização operacional, nunca aceito sem escopo.
\item[Empréstimo] \texttt{id\_almoxarifado}, \texttt{id\_resumo\_reagente}, \texttt{id\_especificacao\_reagente} (strings O), \texttt{densidade\_aplicada} (number positiva C para líquido), \texttt{id\_gestor\_retirada} (string O), \texttt{id\_gestor\_devolucao} (string N até retorno). Snapshot químico da retirada preserva a memória de cálculo; não atualizar retroativamente. Retirante, devolvente e operador são papéis distintos.
\item[Patrimônio] \texttt{nome\_equipamento}, \texttt{letra\_inicial\_nome}, \texttt{predio}, \texttt{andar}, \texttt{sala} (strings O) são projeções de Resumo/Local. \texttt{versao} (inteiro O) detecta conflitos. Histórico embute \texttt{alteracoes} (array de maps campo/valor anterior/novo, admitindo null para ausência) e snapshot do local quando necessário a relatório histórico; não propagar local atual sobre fatos passados.
\item[Requisições] Adição inclui \texttt{justificativa\_resposta} (string N, obrigatória após resposta). Edição inclui \texttt{versao\_bem\_origem} (inteiro O), \texttt{novo\_id\_resumo\_bem\_patrimonial} (string N, reclassificação) e \texttt{novo\_documento\_baixa} (string N, C para proposta de baixa). Projeções nome/plaqueta são identificadores de leitura, não autorização. Resposta nunca inventa foto ausente.
\item[Convite] \texttt{exceder\_capacidade} (boolean O, default false) e \texttt{justificativa\_excecao} (string C quando true), \texttt{aceitado\_por} (string N), \texttt{aceitado\_em} (Timestamp N), \texttt{token\_hash} (string O, SHA-256 do token aleatório). Nunca guardar token em claro. Status aceitado exige UID e instante; matrícula obrigatória na identidade final.
\item[Arquivos] Roteiros, fotos e comprovantes mantêm \texttt{storage\_path} (string O quando upload), \texttt{content\_type} (string O), \texttt{tamanho\_bytes} (inteiro positivo O), \texttt{owner\_uid} (string O), \texttt{criado\_em} (Timestamp O) no registro/metadados aplicável. \texttt{file\_url}/\texttt{photo\_url} são referências legadas; URL assinada transitória não é fonte de autorização nem conteúdo persistente. Arquivo e Firestore não têm transação única: usar upload provisório e reconciliação de órfãos, preservando anexos referenciados.
\item[Etiquetas] \texttt{Contador\_Codigo\_Frasco/singleton}: \texttt{ultimo\_codigo\_gerado} (inteiro não negativo O). Não criar simultaneamente Contadores/codigo\_frasco. Impressao\_Etiqueta\_Frasco acrescenta \texttt{start\_row} (1--10), \texttt{start\_col} (1--3), \texttt{id\_operacao} (string O) e \texttt{status\_geracao} (SOLICITADA/CONCLUIDA/FALHOU). Segunda via fica na auditoria; metadata inclui operação, frasco, motivo e resultado, sem duplicar registro como lote virgem.
\item[Locks] Locks\_Requisicao\_Patrimonio contém \texttt{id\_requisicao}, \texttt{tipo} (ADICAO/EDICAO), \texttt{chave\_recurso} (string O) e \texttt{criado\_em} (Timestamp O). Lock e requisição são atômicos; limpeza só após verificar inexistência de pendência vinculada.
\item[Unicidade técnica] \texttt{Chaves\_Unicas/chaveHash}: \texttt{tipo}, \texttt{id\_recurso} (strings O), \texttt{criado\_em} (Timestamp O). Chave determinística derivada do valor normalizado e domínio (ex: \texttt{Bem\_Patrimonial\_\_\{numero\}} para plaqueta, matrícula, CAS, matéria, local, código turma, lote ou convite pendente); nunca expor e-mail na chave. Criar/liberar na mesma transação do recurso. A implementação exige \textbf{backfill} populando esta coleção para todos os registros preexistentes antes de ativar a verificação transacional.
\item[Deduplicação de Eventos] \texttt{Eventos\_Processados/eventId}: \texttt{processado\_em} (Timestamp O). Usado em transações originadas por triggers (ex: \texttt{onDocumentCreated}) para converter a entrega "at-least-once" do Cloud Functions em processamento "exactly-once".
\item[Operações técnicas] \texttt{Operacoes/id}: \texttt{uid}, \texttt{acao}, \texttt{recurso}, \texttt{hash\_entrada} (strings O), \texttt{status} (PENDENTE/CONCLUIDA/FALHOU), \texttt{resultado} (map N), \texttt{criado\_em}, \texttt{atualizado\_em} (Timestamp O). Chave é vinculada ao ator/ação; repetição com entrada diferente é rejeitada. Suporta idempotência, notificações e reconciliação Auth/Storage. Leitura restrita ao solicitante por resposta mínima; escrita servidor.
\item[Controle de papéis] \texttt{Controle\_Papeis/singleton}: \texttt{chefes\_ativos}, \texttt{gestores\_patrimoniais\_ativos} (inteiros não negativos O), \texttt{versao} (inteiro O). Atualizado na mesma transação das concessões/revogações para impedir retirada concorrente do último responsável.
\item[Materializações] Chaves por lote; data+almoxarifado; data+resumo+almoxarifado; data+local (sentinela GERAL); gestor+ano+mês. Todos os campos de métricas do dicionário são calculados no servidor. Acrescentar \texttt{calculado\_em} (Timestamp O) e \texttt{versao\_calculo} (inteiro O); reprocessar a mesma janela substitui resultado ou deduplica eventos, sem somar duas vezes. Não somar linha GERAL às linhas locais. Métrica legada de volume total usado é alias do consumo ml e não uma parcela adicional.
\end{description}

\paragraph{Consultas e integridade do dicionário.}
A entidade 3FN mantém dependências e referências canônicas; desnormalização é exceção documentada para consulta, não replicação automática de toda FK. Atualização transacional é obrigatória em vínculos, saldos e exclusões lógicas; propagação de nomes pode ser assíncrona, com versão da origem e reconciliação. Nunca aplicar trigger de nome atual a snapshot histórico. Histórico por bem admite collection-group autorizado para relatórios, portanto subcoleção não impede consulta global.

As consultas UI exigem índices versionados em firestore.indexes.json: empréstimo por almoxarifado/status/data prevista e por retirante/data; histórico de frasco por frasco/timestamp e almoxarifado/timestamp; requisição por solicitante/status/feita\_em; turma por professor/status/ano; notificações por papel/lida/emitida\_em; resumo diário por escopo/data. Acrescentar docId como desempate da paginação. Combinações opcionais de filtros precisam de inventário de consultas e validação em emulador/ambiente de integração, não promessa de qualquer combinação arbitrária.

\endgroup
\subsection{Estratégia de Busca Textual no Firestore}
\label{subsec:busca-textual}

\begin{explicacao}[Por que não existe \texttt{LIKE} no Firestore]
O Firestore não oferece busca por substring (sem \texttt{LIKE}, sem \texttt{icontains}). A estratégia adotada é: aplicar \textbf{filtros obrigatórios de igualdade} no Firestore para reduzir drasticamente o conjunto de documentos retornados, e então aplicar \textbf{substring no frontend} sobre os poucos resultados restantes.
\end{explicacao}

\subsubsection{Reagentes}
O usuário deve obrigatoriamente escolher ao menos um dos filtros: estado físico (\texttt{SOLIDO} ou \texttt{LIQUIDO}), natureza química (\texttt{ORGANICO}, \texttt{INORGANICO}, \texttt{ELEMENTO} ou \texttt{HIBRIDO}), tipo de substância (\texttt{PURA} ou \texttt{MISTURA}) e/ou letra inicial (campo \texttt{letra\_inicial} denormalizado). A query Firestore usa esses filtros de igualdade e retorna poucos documentos; o frontend aplica a busca por substring sobre o campo \texttt{nome}.

\subsubsection{Bens Patrimoniais}
Combinação obrigatória de pelo menos um: (Prédio) ou (letra inicial do nome do equipamento) ou (Status). O Firestore filtra com os campos denormalizados \texttt{predio}, \texttt{andar}, \texttt{sala}, \texttt{status} e \texttt{letra\_inicial\_nome}. O frontend aplica substring sobre \texttt{nome\_equipamento} e \texttt{numero\_patrimonio} nos resultados. Bigramas não serão implementados na V1: o filtro de (Prédio + Andar) já reduz os candidatos suficientemente.

\subsubsection{Alunos}
Busca via campo denormalizado \texttt{letra\_inicial} (filtro de igualdade). O frontend aplica substring sobre \texttt{nome} nos resultados. A \textit{range query} por prefixo (\texttt{nome >= ``A''}, \texttt{nome < ``B''}) é alternativa válida, mas não combina facilmente com outros filtros de igualdade em índice composto no Firestore; o campo \texttt{letra\_inicial} é preferível.

\subsubsection{Professores}
Poucos no sistema (dezenas). Usar listagem autorizada e paginada com campos mínimos; filtrar por substring apenas os registros carregados. UI-13 define cache e ciclo de vida.

\begin{regra}[Listeners e cache com escopo]
Aplicar UI-13: listeners limitados às telas visíveis, histórico paginado, memória por padrão e persistência somente em dispositivo confiável. Cache não garante atualização nem ausência de novas leituras faturadas ao reconectar. Encerrar listeners e limpar dados ao sair ou trocar identidade.
\end{regra}

\subsection{Espelhamento Bidirecional da Associação Aluno--Turma}
\begin{itemize}
    \item \textbf{Justificativa NoSQL:} No modelo 3FN, a tabela associativa \texttt{Aluno\_x\_Turma} responde bidirecionalmente às consultas via chave composta. No Firestore, como subcoleções não podem ser filtradas a partir da coleção pai (\texttt{Turma}), e regras de segurança não atuam como filtros de busca, a relação $N:N$ é intencionalmente denormalizada e espelhada em duas estruturas complementares:
    \begin{enumerate}
        \item \textbf{\nolinkurl{Turma/\{turmaId\}/Alunos/\{alunoId\}} (Visão da Turma):} Utilizada pelo professor para gerenciar os matriculados e validar a regra de capacidade máxima (\texttt{qtd\_alunos < capacidade}, atualizado na transação).
        \item \textbf{\nolinkurl{Usuarios/\{uid\}/Turmas/\{turmaId\}} (Visão do Aluno):} Utilizada exclusivamente pelo frontend do Aluno para alimentar o painel lateral de ``Minhas Turmas'' via listener em tempo real (\texttt{onSnapshot}), contendo os metadados essenciais (\texttt{nome\_turma}, \texttt{nome\_materia}, \texttt{ano}, \texttt{semestre}, \texttt{id\_professor}).
    \end{enumerate}
    \item \textbf{Consistência Transacional:} A criação e a remoção desse espelho são executadas de forma estritamente atômica pelo backend nas Cloud Functions. Na matrícula (\texttt{ingressarEmTurmaPorCodigo}), o espelho é criado. Na remoção (\texttt{removerAlunoTurma}), a transação deve abranger: exclusão de \nolinkurl{Turma/\{id\}/Alunos/\{uid\}}, exclusão de \nolinkurl{Usuarios/\{uid\}/Turmas/\{turmaId\}}, gravação no histórico da turma, decremento do contador de vagas e gravação da auditoria.
\end{itemize}
